\documentclass[11pt,a4paper]{article}
\usepackage{gtu-el-yazisi}
\begin{document}
\notbilgi{Alan Teorisi}{Elektrostatik: Coulomb yasası ve elektrik alan}{Ders3.pdf}

\sayfa{1}
\section{Elektrostatik}

\elcizimi[0.6]{sekiller/Ders3-s1-1.png}[1]

\[ \vec e_{12} = \frac{\vec r_{12}}{|\vec r_{12}|}, \qquad |\vec e_{12}| = |\vec e_{21}| = 1, \qquad \vec e_{21} = -\vec e_{12} \]
\[ \vec r_{12} = r_{12}\cdot\vec e_{12}, \qquad \vec r_{21} = r_{21}\,\vec e_{21} = -r_{12}\,\vec e_{12} \]
\[ r_{12} = |\overrightarrow{AB}| = r_{21} \]

\formulkutu{\varepsilon_0 = \frac{1}{36\pi}\,10^{-9}\ \mathrm{F/m}\quad\text{(Boşluğun elektrik sabiti)}}

\[ \vec r_{12} = \overrightarrow{AB} = r_{12}\cdot\vec e_{12} = -\vec r_{21} = -r_{21}\,\vec e_{21} \]
\[ \vec F_{12} = \frac{1}{4\pi\varepsilon_0}\,\frac{q_1 q_2}{(r_{12})^2}\,\vec e_{12} = -\vec F_{21} \]

\noindent\textit{($\vec F_{12}$: $q_1$'in $q_2$'ye uyguladığı kuvvet; \ $\vec F_{21}$: $q_2$'nin $q_1$'e uyguladığı kuvvet)}

\elcizimi[0.45]{sekiller/Ders3-s1-2.png}[1]

\[ \overrightarrow{AB}, \qquad \vec e_{AB} = \frac{\overrightarrow{AB}}{|\overrightarrow{AB}|} \]
\[ \vec e_{BA} = -\vec e_{AB} = -\frac{\overrightarrow{AB}}{|\overrightarrow{AB}|} = \frac{\overrightarrow{BA}}{|\overrightarrow{AB}|} \]

\noindent $q_1,\,q_2 \gtrless 0$\\
$q_1 > 0$ veya $< 0$\\
$q_2 > 0$ veya $< 0$

\noindent $q_1\cdot q_2 > 0$ ise güçler birbirini iter

\noindent $q_1\cdot q_2 < 0$ ise güçler birbirini çeker $\left(\begin{matrix} q_1^{+}\ q_2^{-}\\ q_2^{+}\ q_1^{-} \end{matrix}\right)$

\elcizimi[0.6]{sekiller/Ders3-s1-3.png}[1]

\sayfa{2}
\noindent İkiden fazla yük halinde superposition ilkesi geçerlidir

\elcizimi[0.6]{sekiller/Ders3-s2-1.png}[2]

\[ \vec F_1 = \vec F_{31} + \vec F_{21}, \qquad \vec F_2 = \vec F_{12} + \vec F_{32}, \qquad \vec F_3 = \vec F_{23} + \vec F_{13} \]

\elcizimi[0.4]{sekiller/Ders3-s2-2.png}[2]

\[ \vec F_q = \frac{q}{4\pi\varepsilon_0}\sum_{n=1}^{N}\frac{q_n}{(r_n)^2}\,\vec e_n \]

\noindent Eğer yükler ayrık değil, sürekli bir yoğunlukta dağılmışsa

\noindent hacimsel yük yoğunluğu: $\rho = \dfrac{dQ}{dv}$ \quad $(\mathrm{C/m^3})$

\elcizimi[0.65]{sekiller/Ders3-s2-3.png}[2]

\[ d\vec F_q = \underbrace{\frac{1}{4\pi\varepsilon_0}}_{9\cdot 10^{9}\,\mathrm{N}}\cdot\frac{q\,dQ}{r^2}\,\vec e = \frac{1}{4\pi\varepsilon_0}\cdot\frac{q\,\rho\,dv}{r^2}\,\vec e \]
\[ \Rightarrow\ \vec F_q = \int_v d\vec F_q = \int_v \frac{1}{4\pi\varepsilon_0}\,\frac{q\,\rho\,dv}{r^2}\cdot\vec e \]
\[ \vec F_q = \frac{q}{4\pi\varepsilon_0}\int_v \rho\,\frac{\vec r}{r^3}\,dv \]
\[ \vec F_q = \frac{1}{4\pi\varepsilon_0}\,q\int_v \frac{\rho}{r^3}\cdot\vec r\,dv \]
\[ \rho = \rho_{\text{iç}} + \rho_{\text{dış}} \]
\[ \vec E_q \mathrel{\widehat{=}} \vec F_q \]

\sayfa{3}
\noindent Elektrostatik alan
\[ \vec E = \frac{1}{4\pi\varepsilon_0}\int \frac{\rho\,\vec r}{r^3}\,dv \qquad\qquad 10^{-15}\,\mathrm{m} < r < \infty \]

\elcizimi[0.3]{sekiller/Ders3-s3-1.png}[3]

\[ \vec E = \frac{1}{4\pi\varepsilon_0}\,\frac{q}{(r_{12})^3}\,\vec r_{12} \]
\[ [\vec E] = \kutula{\mathrm{V/m}} = \frac{\mathrm{N}}{\mathrm{C}}, \qquad [q] = \mathrm{C}, \qquad [F] = \mathrm{N}, \qquad [\varepsilon_0] = \mathrm{F/m} \]
\[ [\vec r] = \mathrm{m} \]

\section{Uygulama}

\sekilnotu[3]{Basılı soru (Pr.-2): bir proton ile elektron arasındaki mekanik ve elektrik çekim kuvvetlerinin oranı; aynısı iki elektron için. Basılı metin dijitalleştirilmedi.}

\[ \frac{2e}{3} = k_1, \qquad \frac{-e}{3} = k_2 \]
\noindent Maddesel ise: \ $N\cdot e$, \ $N \in \mathbb{Z}$
\[ p = 2k_1 + k_2 \]
\[ n = 3k_1 + 2k_2 \]

\noindent $\star$ \textbf{ek not}

\noindent $q_1$ durgun ve $q_2$ hareketliyse Coulomb yasası bir yönde geçerlidir. Durgunun hareketliye olanı geçerlidir.

\elcizimi[0.35]{sekiller/Ders3-s3-2.png}[3]

\[ \vec F_{12} = \frac{1}{4\pi\varepsilon_0}\,\frac{q_1 q_2}{r_{12}^{3}}\,\vec r_{12} \]
\[ \vec F_{12} \neq -\vec F_{21} \]

\sayfa{4}
\[ \vec r = r\,\vec e_r \]

\elcizimi[0.65]{sekiller/Ders3-s4-1.png}[4]

\begin{align*}
\varepsilon_0 &= \left(\tfrac{1}{36\pi}\right)\cdot 10^{-9}\ \mathrm{F/m}\\
m_e &= 9{,}1\cdot 10^{-31}\ \mathrm{kg}\\
m_p = m_n &= 1{,}67\cdot 10^{-27}\ \mathrm{kg}\\
e &= 1{,}6\cdot 10^{-19}\ \mathrm{C}\\
k &= 6{,}67\cdot 10^{-11}\ \frac{\mathrm{N\,m^2}}{\mathrm{kg^2}}
\end{align*}

\[ \vec F_{mec} = k\,\frac{m_e\, m_p}{r^2}\,\vec e_r \]
\[ \vec F_{elec} = \frac{1}{4\pi\varepsilon_0}\,\frac{|q_e|\,|q_p|}{r^2}\,\vec e_r = \frac{1}{4\pi\varepsilon_0}\cdot\frac{e^2}{r^3}\,\vec r \]

\[ \begin{aligned} |\vec F_{mec}| &= F_m\\ |\vec F_{elec}| &= F_{elec} \end{aligned} \qquad \frac{F_{mec}}{F_{elec}} = \frac{\dfrac{k\,m_e\,m_p\,|\vec r|}{r^3}}{\dfrac{1}{4\pi\varepsilon_0}\,\dfrac{e^2}{r^3}\,|\vec r|} = 4\pi\varepsilon_0\,k\,\frac{m_e m_p}{e^2} \]

\noindent\textit{(Kesirde pay ve paydadaki $|\vec r|$ ve $r^3$ çarpanları çizilerek sadeleştirilmiş.)}

\[ \Rightarrow\ F_{elec} = \frac{e^2}{m_e\cdot m_p}\cdot\frac{1}{4\pi\varepsilon_0 k}\,F_{mec} \]
\[ F_{elec} = \frac{(1{,}6)^{2}\cdot 10^{-38}}{(9{,}1)(1{,}67)\,10^{-31}\cdot 10^{-27}}\cdot\frac{10^{11}}{4\pi\varepsilon_0\,6{,}67}\,F_{mec} \]
\[ \approx \frac{10^{40}}{4}\,F_{mec} \]

\end{document}
